% ==================================================================
% id:       lim-035
% topic:    גבולות ורציפות של פונקציות
% subtopic: משפט היינה
% solution: full
% ==================================================================
\textbf{הוכיחו או הפריכו כל אחד מהסעיפים הבאים:}
\begin{enumerate}
\item אם $\ds\lim_{x\to2}f(x)=5$ ו-$(a_n)$ סדרה המקיימת $a_n\to2$, אז $f(a_n)\to5$.
\ifsol
\hint{מה דורש משפט היינה מאיברי הסדרה?}
\begin{solbox}
הטענה אינה נכונה. למשל $f(x)=5$ לכל $x\ne2$ ו-$f(2)=0$, והסדרה הקבועה $a_n=2$. אז $\ds\lim_{x\to2}f(x)=5$, אבל $f(a_n)=0\to0$. במשפט היינה דורשים $a_n\ne2$.
\end{solbox}
\fi
\item אם $\ds\lim_{x\to2}f(x)=5$, אז $\ds\lim_{n\to\infty}f\lr{2+\frac1{n^2}}=5$.
\ifsol
\begin{solbox}
הטענה נכונה. הסדרה $a_n=2+\frac1{n^2}$ שואפת ל-$2$, ו-$a_n\ne2$ לכל $n$. לכן לפי היינה $f(a_n)\to5$.
\end{solbox}
\fi
\item אם קיימת סדרה $a_n\to0$, עם $a_n\ne0$ לכל $n$, כך ש-$f(a_n)\to1$, אז $\ds\lim_{x\to0}f(x)=1$.
\ifsol
\begin{solbox}
הטענה אינה נכונה. למשל $f(x)=\floor{x}+1$. עבור $a_n=\frac1n$ מתקיים $f(a_n)=1\to1$, אבל עבור $b_n=-\frac1n$ מתקיים $f(b_n)=0\to0$. לכן לפי היינה הגבול $\ds\lim_{x\to0}f(x)$ אינו קיים.
\end{solbox}
\fi
\item אם $\ds\lim_{n\to\infty}f\lr{\frac1n}=0$ וגם $\ds\lim_{n\to\infty}f\lr{-\frac1n}=0$, אז $\ds\lim_{x\to0}f(x)=0$.
\ifsol
\hint{נסו להגדיר פונקציה ששווה ל-$0$ רק בנקודות $\pm\frac1n$.}
\begin{solbox}
הטענה אינה נכונה. נגדיר $f(x)=0$ אם $x=\frac1n$ או $x=-\frac1n$ עבור $n$ טבעי כלשהו, ו-$f(x)=1$ בכל נקודה אחרת. שתי הסדרות בטענה מתאפסות. אבל עבור $c_n=\frac{1}{n+\frac12}$ מתקיים $c_n\to0$, ו-$c_n$ אינו מהצורה $\frac1m$, ולכן $f(c_n)=1\to1$. לפי היינה הגבול אינו $0$.
\end{solbox}
\fi
\item (*) אם לכל סדרה $a_n\to0$ עם $a_n\ne0$ הסדרה $\lr{f(a_n)}$ מתכנסת, אז קיים הגבול $\ds\lim_{x\to0}f(x)$.
\ifsol
\hint{קחו שתי סדרות כאלה, ובנו מהן סדרה שלישית שהאיברים שלה לקוחים לסירוגין מכל אחת מהן.}
\begin{solbox}
הטענה נכונה. צריך להראות שכל הסדרות האלה נותנות את אותו גבול. תהיינה $a_n\to0$ ו-$b_n\to0$ עם איברים שונים מ-$0$, ונניח $f(a_n)\to L_1$ ו-$f(b_n)\to L_2$. נגדיר את הסדרה
\[c=\lr{a_1,b_1,a_2,b_2,a_3,b_3,\dots}\]
גם $c_n\to0$ ו-$c_n\ne0$, ולכן לפי הנתון $f(c_n)\to L$ עבור $L$ כלשהו. הסדרות $\lr{f(a_n)}$ ו-$\lr{f(b_n)}$ הן תתי-סדרות של $\lr{f(c_n)}$, ולכן גם הן שואפות ל-$L$. כלומר $L_1=L_2=L$. כל הסדרות האלה נותנות את אותו גבול $L$, ולכן לפי היינה $\ds\lim_{x\to0}f(x)=L$.
\end{solbox}
\fi
\end{enumerate}
